% Introduction
\section{Introduction}

Effective concept erasure must do two things at once: remove the target concept and preserve the rest of the representation. This is especially difficult for nonlinear concepts: an edit may defeat one trained probe while leaving behind signal that a freshly trained nonlinear classifier can still recover.
Recent studies have primarily examined linear concept erasure \citep{ravfogel-etal-2020-null, belrose-etal-2023-leace} \ma{cite}, providing strong theoretical guarantees under various assumptions about when a concept has been removed from a representation \ye{why do the theoretical guarantees matter here? remove}. Yet modern neural networks are fundamentally nonlinear, and the concepts encoded in their activations are therefore often nonlinear as well. This leaves nonlinear concept erasure as an important open problem for practical applications, including bias mitigation, fairness, and decision-making systems involving non-subjective concepts.
\ye{i would start with this last sentence: we want to make models better/more fair for these different applications. one way to do this is to erase the concepts from representations. however, concepts are often non-linear, and while recent work has focused on linear concepts, this doesn't work well when concepts are non-linear}\ma{The issue is that we target static activations in our experiment, and dont have experiments on model predictions using those modified activations}


% as emphasized by Elazar and Goldberg's cautionary analysis of post-hoc recovery after adversarial removal \citep{elazar-goldberg-2018-adversarial}. For erasure to be useful, it must therefore achieve two goals simultaneously: high \emph{erasure precision} on the target concept and high \emph{surgicality} on control concepts. \ye{the focus here on our paper will beg the question of why we didn't compare to that paper's method}
% \ye{but also, it's unclear why is this the first thing to mention in the introduction. this is an important detail for the evaluation, but it's not the core idea of this paper}\ye{the first sentence and a half are good. continue from there}\ye{last sentence is repetitive of the first one. remove.}

We approach this problem through a local manifold view of representation space. Hidden states live in a high-dimensional ambient space \ye{this term isn't clear/known. you should explain what you mean by that, and use italics the first time you use it}, but the representations used by natural inputs need not fill that space. Instead, they may concentrate near a much smaller local support. Under this view, the main geometric risk is not simply that an edit is global, but that an ambient-space update can move representations off the natural support where the concept actually lives.

This suggests a different design principle for nonlinear erasure. Rather than editing along the full ambient gradient \ye{again, unclear}, we estimate a local tangent approximation around each sample and retain only the component of the update supported by that local geometry. In practice, we approximate the tangent space with local PCA/SVD \ye{which one is it? are these two options? can we just say matrix decomposition?} on neighborhood secants. This is a first-order approximation rather than a theorem about hidden-state geometry, but it yields a concrete and testable erasure method.

Our contributions are twofold:
\begin{enumerate}
    \item \textbf{A manifold-motivated formulation of nonlinear concept erasure.} We frame nonlinear erasure through a local manifold hypothesis: the relevant question is not only whether representations are low-rank, but whether editing should respect the local support rather than act directly in ambient space. \ye{why is this not part of the second point?}

    \item \textbf{A tangent-constrained nonlinear erasure method.} In Section~\ref{sec:method}, we introduce \TanGCE{}, which operationalizes this view by projecting ambient erasure updates onto locally estimated tangent spaces so that interventions follow the first-order geometry supported by nearby representations. \ye{it's unclear what you method does. you're using all these terms that you don't explain and just sound fancy. the beauty is in simplicity}
    \ye{also, you didn't mention at all the part about using knn of other inputs to help you locate the manifold. it's very implicit in the writing, and you should be explicit about it}
\end{enumerate}

We therefore evaluate with a recoverability-oriented protocol: after erasure, we retrain a fresh nonlinear probe and iterate the intervention to remove residual signal. Across our completed text and vision evaluations \ye{first time you mention those. it sounds like you already mentioned them. you should have a paragraph stating you run extensive experiments and evaluation on a variety of datasets and modalities, as well as models (of diff families and sizes)}, \TanGCE{} is the stronger nonlinear remover than its ambient ablation \ye{unclear what ambient ablation is} and the earlier methods we compare against. The geometry results \ye{what's that?} support the manifold motivation in a bounded form: representations show local low-rank structure, but ambient edits can still leave the natural support. Section~\ref{sec:related} positions the method against prior erasure work, Section~\ref{sec:method} gives the method and modeling view, Section~\ref{sec:evaluation} gives the experimental setup, and Section~\ref{sec:results} evaluates both erasure performance and the geometry motivating local tangent constraints.
